\chapter{Representation Carriers and Derivation Graphs}
\index{representation carrier|textbf}
\index{derivation graph|textbf}
\index{Bloch operator field|textbf}
\index{cell-hopping field|textbf}

\label{ch:representation-carriers-derivation-graphs}

\chapterstatus{foundational mathematical and scientific vocabulary. This
chapter specifies distinctions used by later derivations; it reports no new
calculation, numerical verification, material result, scientific validation, or
uncertainty quantification.}

\section{Learning goals}

After working through this chapter, a reader should be able to

\begin{enumerate}
  \item distinguish a Bloch-operator field from a cell-hopping coefficient
        field by its carrier domain;
  \item identify the state space, ordered basis, geometry, energy reference,
        scalar field, coordinate convention, rank, and source needed to
        interpret a finite representation node;
  \item distinguish an ambient-space role from a retained-space role without
        treating either label as evidence;
  \item explain why gauge conjugation, alignment, Fourier conversion, support
        truncation, and Fourier reconstruction are directed derivation edges
        rather than mutually exclusive node states;
  \item state what ``complete'' means for a declared finite Fourier inventory;
  \item follow a composed reduction chain without conflating projection,
        basis transformation, alignment, Fourier conversion, and truncation;
        and
  \item separate representation metadata, transformation witnesses, numerical
        diagnostics, qualification decisions, and scientific claims.
\end{enumerate}

Chapter~\ref{ch:training-finite-representations} distinguished an abstract
operator from a matrix in a declared basis.
Chapter~\ref{ch:training-periodic-geometry} supplied the direct and reciprocal
coordinate structure used by periodic fields,
Chapter~\ref{ch:training-bloch-finite-differences} constructed finite Bloch
representations, and Chapter~\ref{ch:training-projection-retained-spaces}
distinguished retained-space projection from eliminated-space reduction.
These prerequisites expose one more layer of care: the matrices usually form a
field indexed either by momenta or by lattice translations. Those two index
sets are not interchangeable. Nor is a label such as ``aligned Hamiltonian''
sufficient to identify the space, basis, domain, and operations that produced
the array.

This chapter introduces a typed mathematical vocabulary for the resulting
chain. The word \emph{typed} is used mathematically: a node states what kind of
field it is, while an edge states what kind of transformation relates two
nodes. The chapter does not prescribe a programming language or wire format.
Its purpose is to keep later equations and evidence claims attached to the
objects they actually concern.

\section{Two carrier families}
\index{carrier domain}
\index{momentum domain}
\index{translation domain}

Let
\begin{equation}
  K=(k_0,\ldots,k_{N-1})
  \label{eq:representation-carrier-momentum-domain}
\end{equation}
be an ordered finite momentum domain, and let
\begin{equation}
  \Gamma=(R_0,\ldots,R_{M-1})
  \label{eq:representation-carrier-translation-domain}
\end{equation}
be an ordered finite inventory of lattice translations. The entries in each
inventory are distinct, and a sum over an inventory means that every declared
index position occurs exactly once. Ordering is part of both declarations
because stored field values are correlated positionally with the corresponding
domain. The coordinate convention is also part of the declaration: reduced
reciprocal coordinates, Cartesian reciprocal coordinates, and integer cell
coordinates cannot be exchanged silently.

Let $V$ be a finite-dimensional represented state space. A Bloch-operator field
has the carrier
\begin{equation}
  H_K:K\longrightarrow\operatorname{End}(V),
  \qquad k\longmapsto H_K(k),
  \label{eq:representation-carrier-bloch-field}
\end{equation}
whereas a cell-hopping coefficient field has the carrier
\begin{equation}
  \widehat H_\Gamma:\Gamma\longrightarrow\operatorname{End}(V),
  \qquad R\longmapsto\widehat H_\Gamma(R).
  \label{eq:representation-carrier-hopping-field}
\end{equation}
Even when every value is stored as an $r\times r$ complex matrix, the two
fields have different domains and different interpretations. Equal matrix
shape does not make them values of one carrier type.

\begin{definition}[Bloch-operator representation]
A finite Bloch-operator representation is an ordered momentum domain $K$, an
identified finite represented space $V$ with an ordered basis, and a matrix
value $H_K(k)\in\mathbb C^{r\times r}$ correlated with each $k\in K$. Its
identity additionally states the geometry, energy reference, scalar field,
coordinate convention, and source from which the representation was obtained.
\end{definition}

\begin{definition}[Cell-hopping representation]
A finite cell-hopping representation is an ordered translation inventory
$\Gamma$, the identified represented space and ordered basis in which the
coefficients act, and a matrix value
$\widehat H_\Gamma(R)\in\mathbb C^{r\times r}$ correlated with each
$R\in\Gamma$. Its identity states the same interpretive metadata required for
a Bloch representation.
\end{definition}

These definitions describe finite matrix representations. They do not identify
a finite field with a continuum operator, prove Hermiticity, or assert that the
declared domain is sufficiently refined.

\subsection{Completeness is relative to a declared finite inventory}
\index{complete declared inventory}

For Fourier conversion, declaration of two inventories is not enough. Let
$K=(k_0,\ldots,k_{N-1})$ be nonempty, so $N>0$, and let
$\Gamma_K=(R_0,\ldots,R_{N-1})$ have the same cardinality. A declared phase
pairing defines the characters
\begin{equation}
  \chi_R(k)=e^{+i\langle k,R\rangle}.
  \label{eq:representation-carrier-character-pairing}
\end{equation}

\begin{definition}[Qualified finite Fourier pair]
The ordered inventories $K$ and $\Gamma_K$, together with the declared phase
pairing, form a qualified finite Fourier pair when
\begin{align}
  \frac{1}{N}\sum_{k\in K}
    \overline{\chi_R(k)}\chi_{R'}(k)
  &=\delta_{R,R'},
  &&R,R'\in\Gamma_K,
  \label{eq:representation-carrier-translation-orthogonality}\\
  \frac{1}{N}\sum_{R\in\Gamma_K}
    \chi_R(k)\overline{\chi_R(k')}
  &=\delta_{k,k'},
  &&k,k'\in K.
  \label{eq:representation-carrier-momentum-orthogonality}
\end{align}
The first identity makes coefficient extraction recover every declared
coefficient after synthesis; the second makes complete-inventory synthesis
recover every declared momentum value after extraction.
\end{definition}

The adjective \emph{qualified} is substantive. In an exact derivation the two
orthogonality identities are hypotheses or proved consequences of the finite
character construction. In a computational record, naming two inventories a
Fourier pair does not establish them: their cardinalities, index correspondence,
phase convention, and orthogonality contract still require the declared exact
or numerical qualification route.

A coefficient field on all of a qualified $\Gamma_K$ may be called
\emph{complete for the declared finite inventory}. This phrase means that no
member of that finite transform inventory has been omitted. It does not mean
that all translations of an infinite lattice have been retained, that a
continuum limit has been reached, or that the parent physical model is complete.

A support-truncated field instead retains a declared subset
$S\subseteq\Gamma_K$. The distinction is therefore
\begin{equation}
  \text{complete on }\Gamma_K
  \quad\hbox{versus}\quad
  \text{restricted to }S\subseteq\Gamma_K,
  \label{eq:representation-carrier-complete-versus-truncated}
\end{equation}
not ``exact physics'' versus ``approximate physics.'' The complete finite
transform may already inherit parent-model and discretization errors.

\section{Identity of a representation node}
\index{representation node}
\index{representation identity}

A representation graph contains nodes for fields and directed edges for the
operations relating them. For the present program, the identity of a node must
retain at least

\begin{enumerate}
  \item an identity for the representation node;
  \item its carrier-domain identity and ordering;
  \item the represented state-space identity;
  \item the ordered basis identity;
  \item geometry and coordinate conventions;
  \item scalar-field identity;
  \item energy zero and energy unit;
  \item represented rank;
  \item a family-local role; and
  \item the identity of the source formula, calculation, or artifact.
\end{enumerate}

The list is deliberately redundant with the cautions of
Chapter~\ref{ch:training-finite-representations}. A graph node must make those
cautions operational. For example, two arrays with the same shape and values
but different ordered bases are not the same represented node. Two coefficient
fields with the same matrices but different translation ordering are also not
the same represented node.

\subsection{Family-local roles}
\index{ambient space}
\index{retained space}

Within the Bloch family, the labels \emph{ambient} and \emph{retained} are
useful space roles. An ambient node represents the declared parent fiber space.
A retained node represents a declared retained fiber space. The actual state
space, basis, and rank must still be recorded; the role label does not prove
them.

Within the cell-hopping family, the corresponding useful roles are
\emph{complete declared inventory} and \emph{support truncated}. These roles
refer to coefficient-domain coverage, not to how the coefficient field was
constructed or whether it passed a locality criterion.

The two role families should not be merged into a single list of interchangeable
states. Ambient and retained refer to represented Bloch spaces. Complete and
support truncated refer to translation-domain inventories. Their meanings are
tied to different carriers.

\section{Transformation history belongs to edges}
\index{transformation edge}
\index{derivation history}

A directed transformation edge records
\begin{equation}
  e=
  \bigl(
    \text{source},\text{target},\text{kind},
    \text{method},\text{orientation}
  \bigr),
  \label{eq:representation-carrier-edge-identity}
\end{equation}
together with any witness, diagnostics, and qualification evidence required by
the transformation contract. Source and target are full node identities rather
than bare names. Orientation includes conventions such as whether a gauge acts
as $G^\dagger H G$ or $G H G^\dagger$.

The principal transformation kinds in the periodic reduction chain are

\begin{description}
  \item[Projector construction.] Construction of an ambient-space projector
    $P=FF^\dagger$ from a qualified isometric frame $F$.
  \item[Frame compression.] Passage from an ambient operator to its retained
    coordinate matrix $F^\dagger H F$.
  \item[Invariant restriction.] Identification of a frame compression with the
    restriction of an ambient operator to the frame range, available only when
    that range is invariant under the operator.
  \item[Gauge conjugation.] A unitary change of coordinates within an identified
    retained space.
  \item[Alignment.] A corrective map chosen from a declared admissible family
    and objective.
  \item[Finite Fourier conversion.] A conversion from a Bloch-domain field to
    a cell-hopping field on a declared finite transform inventory.
  \item[Support truncation.] Restriction of the coefficient inventory to a
    declared support.
  \item[Fourier reconstruction.] Construction of a Bloch-domain field from a
    declared coefficient field.
\end{description}

These are edge kinds, not mutually exclusive labels for a node. A retained
Bloch representation may have been gauge conjugated, aligned, and reconstructed.
Encoding those facts as one exclusive node state would discard valid derivation
history. A directed multigraph permits each operation to retain its own source,
target, method, orientation, and evidence.

\begin{remark}[A derivation graph need not be acyclic]
Under the qualified finite Fourier-pair identities, conversion and
reconstruction form an exact inverse pair on the complete declared inventory.
The representation structure is therefore a
directed graph, and potentially a directed multigraph, rather than necessarily
a directed acyclic graph. Repeated methods or alternative routes may also
connect the same pair of nodes.
\end{remark}

\section{The controlled periodic chain}
\index{gauge conjugation}
\index{Fourier conversion}
\index{support truncation}
\index{Fourier reconstruction}

Let $V_A$ and $V_B$ be identified finite-dimensional complex spaces with
$\dim V_A=D$ and $\dim V_B=r$. The ambient and retained node identities record
their ordered basis families $\mathcal B_A(k)$ and $\mathcal B_B(k)$. The
ambient field and retained frame have the explicit types
\begin{equation}
  H_A:K\longrightarrow\operatorname{End}(V_A),
  \qquad
  F(k):V_B\longrightarrow V_A.
  \label{eq:representation-carrier-ambient-frame-types}
\end{equation}
The columns of the $D\times r$ matrix $F(k)$ are the retained ordered basis
vectors expressed in $\mathcal B_A(k)$; coordinate columns therefore obey
$c_A=F(k)c_B$. For a qualified frame edge, $F(k)$ is an isometry,
$F(k)^\dagger F(k)=I_{V_B}$.

Projector construction and frame compression are different operations. The
former creates the ambient-space projector
\begin{equation}
  P(k)=F(k)F(k)^\dagger\in\operatorname{End}(V_A),
  \label{eq:representation-carrier-projector-construction}
\end{equation}
whereas the latter creates the retained reference
\begin{equation}
  H_B(k)=F(k)^\dagger H_A(k)F(k)
  \in\operatorname{End}(V_B).
  \label{eq:representation-carrier-retained-reference}
\end{equation}
Equation~\eqref{eq:representation-carrier-retained-reference} is a frame
compression whether or not the frame range is invariant. It represents an
invariant restriction only under the additional condition
$H_A(k)\operatorname{ran}F(k)\subseteq\operatorname{ran}F(k)$. None of these
operations is a gauge transformation or a support truncation.

The displayed chain below is an exact conditional mathematical chain. A runtime
array supplied as $F$, $G$, or $U$ is only a candidate witness until the required
isometry, unitarity, orientation, domain coverage, and transformation contract
have been qualified by the declared evidence route. Node metadata and role
labels do not supply those facts.

For a qualified gauge-conjugation edge, let $\mathcal B_G(k)$ be another
ordered basis of $V_B$ and let $G(k)$ be pointwise unitary. Its columns are the
vectors of $\mathcal B_G(k)$ expressed in $\mathcal B_B(k)$, so
$c_B=G(k)c_G$. The gauge-conjugated representation is
\begin{equation}
  H_G(k)=G(k)^\dagger H_B(k)G(k)
  \in\operatorname{End}(V_B).
  \label{eq:representation-carrier-gauge-conjugation}
\end{equation}
Both $H_B$ and $H_G$ remain retained Bloch representations. Their different
derivation histories are carried by the gauge-conjugation edge.

For a qualified alignment edge, let $\mathcal B_U(k)$ be the aligned ordered
basis, and let $U(k)$ be pointwise unitary with columns expressing
$\mathcal B_U(k)$ in $\mathcal B_G(k)$, so $c_G=U(k)c_U$. Membership in the
declared alignment family and acceptance under the declared objective,
orientation, and uniqueness policy are separate qualification obligations.
The aligned-coordinate representation is
\begin{equation}
  H_U(k)=U(k)^\dagger H_G(k)U(k)
  \in\operatorname{End}(V_B).
  \label{eq:representation-carrier-alignment}
\end{equation}
Again, $H_U$ remains a retained Bloch representation. Calling it only
``aligned'' would omit the state-space role and would hide the edge on which the
alignment witness and qualification belong.

Assume now the qualified finite Fourier-pair contract above. The basis family
$\mathcal B_U(k)$ supplies one declared ordering of retained coordinate labels
across $K$; Fourier conversion acts entrywise in those identified labels. With
the selected sign and normalization convention, define
\begin{equation}
  \widehat H:\Gamma_K\longrightarrow\operatorname{End}(V_B),
  \qquad
  \widehat H(R)
  =\frac{1}{N}\sum_{k\in K}
    \overline{\chi_R(k)}H_U(k).
  \label{eq:representation-carrier-forward-fourier}
\end{equation}
This edge changes carrier family: its source is a Bloch field and its target is
a cell-hopping field complete on the declared finite inventory. The sign,
normalization, mesh ordering, translation ordering, coordinate pairing, and
retained-label identification are part of the transformation convention.

For $S\subseteq\Gamma_K$, with inherited translation identities and ordering,
define the support-truncated coefficient field as the restriction
\begin{equation}
  \widehat H_S:S\longrightarrow\operatorname{End}(V_B),
  \qquad
  \widehat H_S(R)=\widehat H(R)\quad\text{for }R\in S.
  \label{eq:representation-carrier-support-truncation}
\end{equation}
This edge remains within the cell-hopping family while changing the declared
translation domain from $\Gamma_K$ to $S$. It does not change the retained
state-space rank and is not another projection in state space. An implementation
may equivalently store a zero extension on $\Gamma_K$, but it must declare that
storage convention rather than confuse the zero extension with the support
itself.

The corresponding reconstructed Bloch field has type
$H_C(U,S):K\to\operatorname{End}(V_B)$ and is
\begin{equation}
  H_C(U,S;k)
  =\sum_{R\in S}\chi_R(k)\widehat H_S(R).
  \label{eq:representation-carrier-reconstruction}
\end{equation}
Fourier reconstruction changes the carrier back to a Bloch-domain field. The
target remains a retained Bloch representation; reconstruction is recorded by
the incoming edge rather than by replacing its space role. On the complete
qualified inventory, the momentum orthogonality identity gives the exact
finite reconstruction
\begin{equation}
  H_C(U,\Gamma_K;k)
  =\frac{1}{N}\sum_{k'\in K}
    \left(\sum_{R\in\Gamma_K}
      \chi_R(k)\overline{\chi_R(k')}\right)H_U(k')
  =H_U(k).
  \label{eq:representation-carrier-complete-reconstruction}
\end{equation}
No corresponding equality is asserted for a proper support $S$ without an
additional support or residual hypothesis.

The full typed chain is therefore
\begin{equation}
  \begin{aligned}
    H_A
    &\xrightarrow{\text{frame compression}} H_B
     \xrightarrow{\text{gauge conjugation}} H_G
     \xrightarrow{\text{alignment}} H_U,\\
    H_U
    &\xrightarrow{\text{finite Fourier conversion}} \widehat H
     \xrightarrow{\text{support truncation}} \widehat H_S
     \xrightarrow{\text{Fourier reconstruction}} H_C(U,S).
  \end{aligned}
  \label{eq:representation-carrier-full-chain}
\end{equation}
The first arrow uses the declared frame $F$ and must not be replaced by the
separate projector-construction edge. Every subsequent comparison must identify
which nodes are being compared and in which common representation.

\section{What the graph does not prove}
\index{evidence boundary!representation graph}

A well-formed identity or edge is necessary for interpretation, but it is not
scientific evidence. In particular, graph construction alone does not prove

\begin{itemize}
  \item Hermiticity, idempotence, orthogonality, unitarity, or determinant
        constraints;
  \item correctness or uniqueness of an alignment solution;
  \item Fourier index bijection, inversion closure, or reconstruction accuracy;
  \item support completeness under a physical metric or symmetry;
  \item numerical convergence or independence of withheld evaluation samples;
  \item agreement with a continuum parent or a material calculation; or
  \item scientific validity or uncertainty bounds.
\end{itemize}

Those statements require separate property evidence with declared tolerances,
methods, and dispositions. A useful evidence sequence is
\begin{equation}
  \text{identity}
  \longrightarrow
  \text{constructed value}
  \longrightarrow
  \text{diagnostics}
  \longrightarrow
  \text{qualification}
  \longrightarrow
  \text{bounded claim}.
  \label{eq:representation-carrier-evidence-sequence}
\end{equation}
No arrow may be inferred merely from the presence of the object on its left.
If a mandatory transformation is rejected or unresolved, it supplies no
qualified target for later Fourier conversion, support truncation,
reconstruction, or model comparison.

\section{Consequences for the reductions that follow}

The carrier distinction organizes the remaining reduction chapters.
Chapter~\ref{ch:training-isolated-band-fourier-reduction} specializes the
Fourier edges to a scalar isolated band, and
Chapter~\ref{ch:training-composite-band-gauge-alignment} develops the
higher-rank frame, gauge, alignment, and block-hopping structure.

Later comparisons inherit the same discipline. Before subtracting two fields,
one must align their carrier domains, represented spaces, ordered bases,
geometries, energy references, units, and coordinate conventions. Before
claiming locality, one must say whether the claim concerns a complete declared
coefficient inventory or a support-truncated one. Before attributing a result
to alignment or truncation, one must retain the separate edges so that the two
operations cannot be silently conflated.
